\section{数字化基础与关键技术}
\subsection{数字化手段与关键装备}
要将动力电池包装配从局部自动化提升到智能装配，首先需要建立稳定的数字化底座。该底座并不只对应某一台机器人或某一套视觉系统，而是覆盖工艺数据采集、任务调度、检测反馈、追溯管理和异常处置的完整链路。典型手段包括条码或 RFID 绑定、MES 工单与配方管理、工位 PLC 与边缘控制器、力矩枪与焊接设备数据采集、视觉检测系统、工业机器人、协作机器人、AGV/AMR 物流单元以及质量数据库等 \cite{Ayerbe2021AEM,Tao2020Sustainability,Park2023RSER,Liu2020Access}。

如果将产线拆分来看，数字化装备大致可以分为三类。第一类是执行装备，如工业机器人、协作机器人、焊接设备、力矩工具和输送单元，主要负责稳定完成工艺动作。第二类是感知装备，如工业相机、激光位移传感器、力传感器和电参数采集设备，主要负责采集工位状态和质量状态。第三类是管理与控制装备，如 MES、PLC、边缘控制器和追溯数据库，主要负责贯通工位之间的数据。这样的分类方式更便于说明后续技术应用所依托的基础条件。

在动力电池包装配现场，关键不在于设备数量本身，而在于设备组合是否合理。例如，模组入箱工位如果只有机器人而缺少定位治具和视觉校正，重复精度便很难真正转化为装配精度；焊接工位如果只有焊机而缺少电阻检测和焊缝视觉检查，也难以形成可靠闭环。对于动力电池包而言，较为典型的装备组合是“机器人/协作机器人 + 工装夹具 + 视觉检测 + 工艺数据库”。这类组合能够同时覆盖定位、执行、检测和追溯四个环节，因此更接近智能装配而非单点自动化。

数字化手段的关键不在于传感器数量，而在于不同设备之间能否形成可追溯的数据链。若焊接设备、力矩设备、视觉设备和装配工位的结果不能绑定到统一产品 ID 上，那么后续异常分析和工艺优化便会失去依据。因此，本节的重点并非罗列设备名称，而是说明动力电池包装配为何能够被持续感知、准确记录、有效分析并形成闭环优化 \cite{Yao2023JMSY}。只有明确这一基础，后文关于 AI、知识图谱和大模型智能体的应用才具有系统支撑。

\begin{figure}[H]
  \centering
  \includegraphics[width=0.70\textwidth]{fig14_plug_produce_cell.png}
  \caption{模块化 Plug \& Produce 制造单元示意}
  \label{fig:plug-produce-cell}
\end{figure}

当产线需要兼顾换型与扩展时，单个设备能力之外还必须关注单元级组织方式。图~\ref{fig:plug-produce-cell}所示的 Plug \& Produce 单元体现了一种较有代表性的模块化布局：围绕共享机器人布置工位、缓存区和物料源，并通过上层规划与调度逻辑快速接入新任务或新资源。这类结构对于动力电池包装配具有启发意义，因为模组入箱、紧固、检测和返修等任务往往并非固定不变，采用可重构单元有助于在不完全推翻原有产线的前提下实现局部升级与柔性扩展 \cite{Massouh2025JIM,Yao2023JMSY,Tao2020Sustainability}。

图~\ref{fig:equipment-placeholder}用于展示“机器人 + 视觉 + 检测”一体化装备单元的典型形式。图中左侧体现的是检测点选择与路径规划，右侧则对应实体机器人在工件表面执行无损检测任务。也就是说，这类单元并不是单纯把相机装在机器人上，而是将离线规划、在线感知和检测执行整合为一个连续过程。图中案例直接来自视觉引导的机器人无损检测研究 \cite{Khan2021JRem}；从更宽泛的制造场景看，这类单元同样可以推广到面向子装配体的自动检查任务 \cite{Ronai2023SAS}，能够直观对应本节讨论的执行装备、感知装备与检测反馈关系。\vspace{-0.4em}

\begin{figure}[H]
  \centering
  \includegraphics[width=0.50\textwidth]{fig3_robotic_ndt_inspection.png}
  \caption{检测点规划与视觉引导机器人检测单元}
  \label{fig:equipment-placeholder}
\end{figure}
